\begin{textsample}{POS Dim 1 – vem\_ed\_05 – Score 5.00 – vem\_ed\_05/t017.txt}  \label{ex:f1_pos_105}
QUEBRA-GELO PROCURA-SE Quer desenvolver um \textbf{intercâmbio} virtual?
Instituições de Ensino estrangeiras buscam professores nas Fatecs em diversas áreas, como: \textbf{Competências} interculturais Mineração e geoprocessamento Mapa de riscos Saúde Artes Álgebra e Cálculo História e Ciência Política Caso tenha interesse, basta escrever para cesu.pci@cps.sp.gov.br e juntar- se às nossas equipes no Teams.
QUEBRA-GELO Começamos 2021 com novidades.
A \textbf{partir} desta edição, oficializamos o nome desta publicação como a chamamos, carinhosamente, na equipe dos Projetos Colaborativos Internacionais: VEm – uma alusão às iniciais de Virtual Exchange, Intercâmbios Virtuais, e “m” para meio de comunicação.
Uma de nossas metas para este ano \textbf{é} \textbf{expandir} ainda mais \textbf{ações} em países hispanofalantes.
Por isso, a seção"Quem \textbf{é} Quem"apresenta Brenda Garcia Portillo, da Universidad de Monterrey (UDEM), instituição mexicana \textbf{que} \textbf{é} \textbf{referência} mundial em Intercâmbios Virtuais.
Em"Boas Práticas", Joelson Alves do Nascimento (Fatec Barueri) relata o projeto sobre \textbf{competências} emocionais com a universidade colombiana Uniminuto.
Sua receita para o sucesso de um \textbf{intercâmbio} virtual soma planejamento e paixão pelo tema a \textbf{ser} pesquisado.
Em dezembro de 2020, realizamos uma pesquisa de percepção com professores e estudantes envolvidos nos Projetos Colaborativos Internacionais desenvolvidos no segundo semestre.
Responderam 25 dos 34 docentes participantes e quase 400 dos 833 alunos envolvidos.
Os principais resultados estão resumidos na página 3.
No último mês de um ano, ao mesmo tempo desafiador e repleto de realizações, perdemos o professor Antal György Almásy, da Fatec Diadema.
Ele \textbf{foi} presidente da Associação Brasileira de Aerossóis e Saneantes Domissanitários (ABAS) e tinha esboçado um projeto com a Durban University of Technology (DUT), na África do Sul, para \textbf{produção} de sabonetes artesanais por comunidades carentes, valorizando a biodiversidade \textbf{local}.
A ele, prestamos nossa homenagem e agradecimento.

% matched lemmas: ação, competência, expandir, intercâmbio, local, partir, produção, que, referência, ser
\end{textsample}
